\chapter{Respuesta en frecuencia de sistemas LIT TC}

\section{Respuesta en frecuencia}

Si la entrada a un sistema de tiempo continuo es $x\left( t\right) =e^{st},$ 
entonces la salida es $y\left( t\right) =H\left( s\right) e^{st}$ donde
\begin{equation}
        H\left( s\right) =\int_{-\infty }^{\infty }h\left( \tau \right) e^{-s \, \tau} \, d\tau ,
\end{equation}
en la cual $h\left( t\right) $ es la respuesta al impulso del sistema LIT.
En el caso general de $s$ complejo, $H\left( s\right) $ es la funci\'{o}n del sistema. 
En el caso $\Re\left( s\right) =0\rightarrow s=j \, \omega $ y $x\left( t\right) =e^{s \, t}=e^{j \, \omega \, t},$ 
la respuesta en frecuencia del sistema esta definida por 
\begin{equation}
        H\left( j\omega \right) =\int_{-\infty }^{\infty }h\left( t\right) e^{-j\omega t}dt.
\end{equation}
La salida $y\left( t\right) $ para una entrada 
\begin{equation}
        x\left( t\right) =\sum_{k=-\infty }^{\infty }a_{k}e^{jk\omega t},
\end{equation}
est\'{a} dada entonces por 
\begin{equation}
        y\left( t\right) =\sum_{k=-\infty }^{\infty }a_{k}H\left( j\omega k\right) e^{jk\omega t}.
\end{equation}

\subsection{Ejemplo}

\begin{enumerate}
\item Sea un sistema con respuesta al impulso dada por 
\begin{equation}
h\left( t\right) =e^{-t}u\left( t\right) ,
\end{equation}
y una entrada dada por 
\begin{equation}
x\left( t\right) =\sum_{k=-3}^{3}a_{k}e^{jk2\pi t},
\end{equation}
donde 
\begin{equation}
\begin{array}{ccccc}
a_{0} &=&1,&& \\
&&&& \\
a_{1} &=&a_{-1}&=&\frac{1}{4}, \\
&&&& \\
a_{2} &=&a_{-2}&=&\frac{1}{2}, \\
&&&& \\
a_{3} &=&a_{-3}&=&\frac{1}{3}.
\end{array}
\end{equation}

\item Calculamos primero la respuesta en frecuencia planteando la integral
\begin{equation}
H\left( j\omega \right) =
  \int_{-\infty }^{\infty }h\left( t\right)e^{-j\omega t}dt.
\end{equation}
El sistema dado es causal porque la respuesta al impulso $\delta(t)$ es nula para $t < 0$. Por lo
tanto, los l\'{i}mites de la integral camb\'{\i}an y resulta
\begin{equation}
H\left( j\omega \right) =
  \int_{0}^{\infty }e^{-\tau }e^{-j\omega \tau}d\tau.
\end{equation}
Resolviendo la integral podemos escribir
\begin{equation}
H\left( j\omega \right) =
  -\frac{1}{1+j\omega }\left[ e^{-\left( 1+j\omega\right) \tau }\right] _{0}^{\infty }.
\end{equation}
Aplicando la regla de Barrow, obtenemos finalmente
\begin{equation}
H\left( j\omega \right) =\frac{1}{1+j\omega }.
\end{equation}

\item Entonces la salida $y\left( t\right) $ es 
\begin{equation}
y\left( t\right) = \sum_{k=-3}^{3}b_{k}e^{jk2\pi t},
\end{equation}
donde, en general,
\begin{equation}
b_{k} = a_{k} H\left( jk2\pi \right).
\end{equation}
Por lo tanto, calculando cada $b_{k}$ obtenemos
\begin{equation}
b_{0} = 1,
\end{equation}
y
\begin{equation}
  \begin{array}{ccc}
  b_{1} =\frac{1}{4}\left( \frac{1}{1+j2\pi }\right), &&
  b_{-1}=\frac{1}{4}\left( \frac{1}{1-j2\pi }\right), \\
  \\
  b_{2} =\frac{1}{2}\left( \frac{1}{1+j4\pi }\right), &&
  b_{-2}=\frac{1}{2}\left( \frac{1}{1-j4\pi }\right), \\
  \\
  b_{3} =\frac{1}{3}\left( \frac{1}{1+j6\pi }\right), &&
  b_{-3}=\frac{1}{3}\left( \frac{1}{1-j6\pi }\right).
  \end{array}
\end{equation}

\end{enumerate}

\section{Representaci\'{o}n gr\'{a}fica de la respuesta en frecuencia de un sistema LIT}

La respuesta en frecuencia es una funci\'{o}n compleja y debe ser representada gr\'{a}ficamente en
\begin{enumerate}
        \item   Forma polar: magnitud y fase.
        
        \item   Forma cartesiana: parte real y parte imaginaria.
\end{enumerate}

Ambas representaciones permiten visualizar distintas propiedades de los sistemas LIT. Por tal motivo las
estudiaremos por separado.

\subsection{Definiciones}

\subsubsection{Ganancia o atenuaci\'{o}n}

La ganancia y la atenuaci\'{o}n de un filtro son conceptos complementarios, y sus valores est\'{a}n 
determinados por la magnitud de la respuesta en frecuencia. Dado que la salida de un filtro $y$ para
una entrada peri\'{o}dica $x$ queda determinada por 
\begin{equation}
        y(t) = H(j \omega) x(t),
\end{equation}
para tiempo continuo. El m\'{o}dulo de la funci\'{o}n compleja $\| H \|$ determina la 
variaci\'{o}n del tama\~{n}o de la salida respecto a la entrada. Hablamos de ganancia cuando deseamos 
saber cuanto queda ampliada la salida respecto a la entrada, y hablamos de atenuaci\'{o}n cuando deseamos 
saber cuanto queda disminuida la salida respecto a la entrada. Dado que por lo general
los filtros amplian determinadas frecuencias y disminuyen otras, hablar de ganancia o atenuaci\'{o}n es 
equivalente a elegir un punto de vista ($G = 1 - A$).

\subsubsection{Fase}

La fase de un filtro es la fase de la respuesta en frecuencia (funci\'{o}n compleja) $H(j \omega)$ en 
tiempo continuo o $H(e^{j \omega})$ para tiempo discreto.

\subsubsection{Fase lineal y no lineal}

Cuando la fase del filtro es una funci\'{o}n \textbf{lineal} de $\omega$, hay una interpretaci\'{o}n 
simple de su efecto en el dominio del tiempo. Consideremos el sistema de tiempo continuo con respuesta en 
frecuencia
\begin{equation}
        H( j \omega) = e^{j \omega t_{0}},
\end{equation}
donde la ganancia es $1$ para toda frecuencia y fase lineal
\begin{equation}
        |H(j \omega)| = 1, \angle H(j w) = -\omega t_{0}.
\end{equation}
La salida de este sistema es igual a su entrada desplazada en el tiempo
\begin{equation}
        y(t) = x(t - t_{0}.
\end{equation}
% ver ejemplo 4.15 en el oppenheim, willsky y nawab

\subsubsection{Retraso de grupo}
\begin{equation}
        \tau (\omega) = -\frac{d \angle H(j \omega)}{d\omega}.
\end{equation}

\section{Gr\'{a}ficos de Bode}

\subsection{Sistemas de primer orden de tiempo continuo}

Consideremos la siguiente ecuaci\'{o}n diferencial 
\begin{equation}
\tau \frac{dy\left( t\right) }{dt}+y\left( t\right) =x\left( t\right) .
\end{equation}

\begin{enumerate}
\item  La funci\'{o}n de transferencia es 
\begin{equation}
        H\left( s\right) =\frac{1}{s\tau +1},
\end{equation}
con un polo en 
\begin{equation}
        s=-\frac{1}{\tau }.
\end{equation}

\item  La funci\'{o}n de respuesta en frecuencia es 
\begin{equation}
H\left( j\omega \right) =\frac{1}{j\omega \tau +1}.
\end{equation}

\item  La respuesta al impulso es 
\begin{equation}
        h_{1}\left( t\right) =\frac{1}{\tau }e^{-\frac{t}{\tau }}u\left( t\right) .
\end{equation}

\item  La respuesta al escal\'{o}n es 
\begin{equation}
        h_{2}\left( t\right) =h_{1}\left( t\right)  u\left( t\right) =
                \left(1-e^{-\frac{t}{\tau }}\right) u\left( t\right) .
\end{equation}

\item  El parametro $\tau $ es la constante de tiempo del sistema y controla la derivada de la salida. 
Comparemos la respuesta al impulso con la respuesta en frecuencia 
\begin{eqnarray}
        \tau h\left( t\right) &=&e^{-\frac{t}{\tau }}u\left( t\right) , \\
        H\left( j\omega \right) &=&\frac{1}{j\omega \tau +1}.
\end{eqnarray}
\end{enumerate}

\paragraph{La funci\'{o}n magnitud de Bode para sistemas de primer orden}
Consideremos la funci\'{o}n magnitud de Bode 
\begin{equation}
        20\log _{10}\left\| H\left( j\omega \right) \right\| =
                -10\log _{10}\left(\left( \omega \tau \right) ^{2}+1\right) .  \label{magbode}
\end{equation}

\paragraph{Simplificaci\'{o}n de la funci\'{o}n magnitud de Bode para sistemas de primer orden}
La funci\'{o}n (\ref{magbode}) puede simplificarse seg\'{u}n el valor que toma la variable $\omega $
\begin{enumerate}
   \item  (forma normalizada) 
   \begin{equation}
      \omega \tau <1\rightarrow 20\log _{10}\left\| H\left( j\omega \right)\right\| 
          \approx -10\log _{10}\left( 1\right) =0,
   \end{equation}
   o (forma no normalizada) 
   \begin{equation}
      \omega <\frac{1}{\tau }\rightarrow 20\log _{10}\left\| H\left( j\omega\right) \right\| \approx 0.
   \end{equation}

   \item  (forma normalizada) 
   \begin{equation}
      \omega \tau >1\rightarrow 20\log _{10}\left\| H\left( j\omega \right)\right\| 
          \approx -20\log _{10}\left( \omega \tau \right) ,
   \end{equation}
   o (forma no normalizada) 
   \begin{equation}
      \omega >\frac{1}{\tau }\rightarrow 20\log _{10}\left\| H\left( j\omega\right) \right\| 
          \approx -20\log _{10}\left( \omega \right) + 20\log_{10}\left( \frac{1}{\tau }\right) .
   \end{equation}
\end{enumerate}
Esto significa que la funci\'{o}n (\ref{magbode}) tiene dos as\'{\i}ntotas. 
Dado que el gr\'{a}fico se realiza con escalas logaritmicas en ambos ejes, estas as\'{\i}ntotas son rectas.

\paragraph{As\'{\i}ntotas de la funci\'{o}n magnitud de Bode para sistemas de primer orden}
La funci\'{o}n (\ref{magbode}) puede aproximarse por as\'{\i}ntotas seg\'{u}n el valor que toma la 
variable $\omega.$
\begin{enumerate}
   \item  As\'{\i}ntota de baja frecuencia, coincidente con la l\'{\i}nea de $0db, $
   \item  As\'{\i}ntota de alta frecuencia, con una pendiente de $20db$ por d\'{e}cada.
   \item  Estas as\'{\i}ntotas se intersectan cuando $\omega \tau =1$ (forma normalizada) o 
     $\omega =\frac{1}{\tau }$ (forma no normalizada).
   \item  El punto de intersecci\'{o}n de las as\'{\i}ntotas se llama frecuencia de quiebre.
\end{enumerate}

Estas as\'{\i}ntotas nos brindan la posibilidad de aproximar r\'{a}pidamente la 
funci\'{o}n (\ref{magbode}). En la frecuencia de quiebre, tenemos adem\'{a}s que el verdadero valor de la 
funci\'{o}n (\ref{magbode}) es 
\begin{equation}
    20\log _{10}\left\| H\left( j\omega \right) \right\| =-10\log _{10}\left(2\right) \approx -3db.
\end{equation}

Consideremos ahora la funci\'{o}n fase de Bode 
\begin{equation}
\angle H\left( j\omega \right) =\arctan \left( \omega \tau \right) . \label{faseBode}
\end{equation}
Esta funci\'{o}n tambien puede aproximarse por distintas rectas para distintos rangos de frecuencia. 
Recordemos el uso de escala logaritmica en el eje de la frecuencia.
\begin{equation}
\begin{array}{r}
   \omega \tau \leq \frac{1}{10}\rightarrow \arctan \left( \omega \tau \right) =0, \\ 
    \frac{1}{10}\leq \omega \tau \leq 10\rightarrow \arctan \left( \omega \tau \right) =
        -\left( \frac{\pi }{4}\right) \left( \log _{10}\left( \omega \tau \right) +1\right) , \\ 
   \arctan \left( \omega \tau \right) =-\frac{\pi }{2}.
\end{array}
\end{equation}
En la frecuencia de quiebre, tenemos adem\'{a}s que el verdadero valor de la funci\'{o}n 
(\ref{faseBode}) es 
\begin{equation}
\angle H\left( j\omega \right) =\arctan \left( 1\right) =\frac{\pi }{4}.
\end{equation}

\subsection{Sistemas de segundo orden de tiempo continuo}

Consideremos la ecuaci\'{o}n diferencial 
\begin{equation}
   \frac{d^{2}y\left( t\right) }{dt^{2}}+
   2\zeta \omega_{n}\frac{dy\left(t\right) }{dt}+
   \omega_{n}^{2}y\left( t\right) =\omega _{n}^{2}x\left(t\right) .
\end{equation}

\begin{enumerate}
\item  La funci\'{o}n de transferencia correspondiente es 
\begin{equation}
  H\left( s\right) =\frac{\omega _{n}^{2}}{s^{2}+2\zeta \omega _{n}s+\omega_{n}^{2}}.
\end{equation}

\item  La respuesta en frecuencia es 
\begin{equation}
   H\left( j\omega \right) =
   \frac{\omega _{n}^{2}}{\left( \omega_{n}^{2}-\omega^{2}\right) +j2\zeta \omega_{n}\omega },
\end{equation}
que se puede factorear como 
\begin{equation}
  H\left( j\omega \right) =
  \frac{\omega _{n}^{2}}{\left( j\omega-c_{1}\right) \left( j\omega -c_{2}\right) },
\end{equation}
donde 
\begin{equation}
c_{1}=\omega _{n}\left( -\xi +\sqrt{\xi ^{2}-1}\right) ,
\end{equation}
y 
\begin{equation}
c_{1}=\omega _{n}\left( -\xi -\sqrt{\xi ^{2}-1}\right) .
\end{equation}

\item La respuesta al impulso la calculamos aplicando expansi\'{o}n en fracciones parciales y la 
antitransformada de Laplace.
\begin{enumerate}
\item  Observemos que 
\begin{equation}
\xi \neq 1\rightarrow c_{1}\neq c_{2},
\end{equation}
y podemos hacer una expansi\'{o}n en fracciones parciales de la forma 
\begin{equation}
H\left( j\omega \right) =\frac{M}{j\omega -c_{1}}-\frac{M}{j\omega -c_{2}},
\end{equation}
donde 
\begin{equation}
M=\frac{\omega _{n}}{2\sqrt{\xi ^{2}-1}}.
\end{equation}
Por lo tanto, la respuesta al impulso para este caso es 
\begin{equation}
h_{1}\left( t\right) =M\left( e^{c_{1}t}-e^{c_{2}t}\right) u\left( t\right) .
\end{equation}

\item  En cambio, si
\begin{equation}
\xi =1\rightarrow c_{1}=c_{2}=-\omega _{n},
\end{equation}
y 
\begin{equation}
H\left( j\omega \right) =
  \frac{\omega _{n}^{2}}{\left( j\omega -\omega_{n}\right) ^{2}},
\end{equation}
y la respuesta al impulso correspondiente es 
\begin{equation}
h_{1}\left( t\right) =\omega _{n}^{2}te^{-\omega _{n}t}u\left( t\right) .
\end{equation}
\end{enumerate}

\item  La respuesta al escal\'{o}n $h_{2}\left( t\right) $ puede calcularse mediante la convoluci\'{o}n 
\begin{equation}
h_{2}\left( t\right) = h_{1}\left( t\right) * u\left( t\right) .
\end{equation}

\begin{enumerate}
\item  En el caso 
\begin{equation}
\xi \neq 1\rightarrow c_{1}\neq c_{2},
\end{equation}
tenemos que la respuesta al impulso es 
\begin{equation}
h_{1}\left( t\right) =M\left( e^{c_{1}t}-e^{c_{2}t}\right) u\left( t\right) ,
\end{equation}
y por lo tanto la respuesta al escal\'{o}n es 
\begin{equation}
h_{2}\left( t\right) =\left( 1+M\left( \frac{e^{c_{1}t}}{c_{1}}-\frac{%
e^{c_{2}t}}{c_{2}}\right) \right) u\left( t\right) .
\end{equation}

\item  En el caso 
\begin{equation}
\xi =1\rightarrow c_{1}=c_{2}=-\omega _{n},
\end{equation}
tenemos que la respuesta al impulso correspondiente es 
\begin{equation}
h_{1}\left( t\right) =\omega _{n}^{2}te^{-\omega _{n}t}u\left( t\right) ,
\end{equation}
y por lo tanto la respuesta al escal\'{o}n es 
\begin{equation}
h_{2}\left( t\right) =\left( 1-e^{-\omega _{n}t}-\omega _{n}te^{-\omega
_{n}t}\right) u\left( t\right) .
\end{equation}
\end{enumerate}

\item  Es posible obtener una expresi\'{o}n normalizada para la respuesta en frecuencia.

\begin{enumerate}
\item  Observe que podemos escribir, dividiendo el numerador y el
denominador por $\omega _{n}^{2},$
\begin{equation}
H\left( j\omega \right) =\frac{1}{\left( \frac{j\omega }{\omega _{n}}%
\right) ^{2}+2\xi \left( \frac{j\omega }{\omega _{n}}\right) +1},
\end{equation}
obteniendo la forma normalizada de un sistema de segundo orden.

\item  En este caso, la respuesta en frecuencia es funci\'{o}n de $\frac{\omega }{\omega _{n}}$. 
El cambio de $\omega _{n}$ equivale a cambiar una escala en tiempo y frecuencia.
\end{enumerate}

\item  El par\'{a}metro $\xi $ es el coeficiente de amortiguamiento, y $\omega _{n}$ es la 
frecuencia natural de oscilaci\'{o}n. La raz\'{o}n de estos nombres se puede entender observando las 
respuestas al impulso y al escal\'{o}n. 

\begin{enumerate}
\item  En el caso $0<\xi <1,$ los coeficientes $c_{1}$ y $c_{2}$ son
complejos, y la respuesta al impulso puede escribirse en la forma
\begin{equation}
h_{1}\left( t\right) =\frac{\omega _{n}e^{-\xi \omega _{n}t}}{2j\sqrt{1-\xi
^{2}}}\left( e^{j\omega _{n}\sqrt{1-\xi ^{2}}t}-e^{j\omega _{n}\sqrt{1-\xi
^{2}}t}\right) u\left( t\right) .
\end{equation}
Reacomodando la expresi\'{o}n, obtenemos
\begin{equation}
h_{1}\left( t\right) =\frac{\omega _{n}e^{-\xi \omega _{n}t}}{\sqrt{1-\xi
^{2}}}\frac{\left( e^{j\omega _{n}\sqrt{1-\xi ^{2}}t}-e^{j\omega _{n}\sqrt{%
1-\xi ^{2}}t}\right) }{2j}u\left( t\right) ,
\end{equation}
equivalente a
\begin{equation}
h_{1}\left( t\right) =\frac{\omega _{n}}{\sqrt{1-\xi ^{2}}}e^{-\xi \omega
_{n}t}\sin \left( \omega _{n}\sqrt{1-\xi ^{2}}t\right) u\left( t\right) .
\end{equation}
Esto significa que el movimiento es oscilatorio amortiguado, y se dice en
este caso que el sistema es subamortiguado.

\item  En el caso $\xi >1,$ los coeficientes $c_{1}$ y $c_{2}$ son negativos
y la respuesta al impulso es la diferencia entre dos exponenciales
decrecientes 
\begin{equation}
h_{1}\left( t\right) =M\left( e^{-\left| c_{1}\right| t}-e^{-\left|
c_{1}\right| t}\right) u\left( t\right) .
\end{equation}
En este caso el sistema es sobreamortiguado.

\item  En el caso $\xi =1,$ cuando $c_{1}=c_{2}=-\omega _{n},$ decimos que
el sistema es criticamente amortiguado 
\begin{equation}
h_{1}\left( t\right) =\omega _{n}^{2}te^{-\omega _{n}t}u\left( t\right) .
\end{equation}
\end{enumerate}

\item  La funci\'{o}n magnitud de Bode para un sistema con respuesta en
frecuencia normalizada  
\begin{equation}
H\left( j\omega \right) =\frac{1}{\left( 1-\left( \frac{\omega }{\omega
_{n}}\right) ^{2}\right) +j2\zeta \left( \frac{\omega }{\omega _{n}}\right) 
},
\end{equation}
est\'{a} dada por
\begin{equation}
20\log _{10}\left\| H\left( j\omega \right) \right\| =-20\log _{10}\left[ 
\sqrt{\left( 1-\left( \frac{\omega }{\omega _{n}}\right) ^{2}\right)
^{2}+4\zeta ^{2}\left( \frac{\omega }{\omega _{n}}\right) ^{2}}\right] ,
\end{equation}
que se simplifica
\begin{equation}
20\log _{10}\left\| H\left( j\omega \right) \right\| =-10\log _{10}\left[
\left( 1-\left( \frac{\omega }{\omega _{n}}\right) ^{2}\right) ^{2}+4\zeta
^{2}\left( \frac{\omega }{\omega _{n}}\right) ^{2}\right] .
\end{equation}
Note que las as\'{\i}ntotas para esta funci\'{o}n son
\begin{equation}
\omega \ll \omega _{n}\rightarrow 20\log _{10}\left\| H\left( j\omega
\right) \right\| \approx 0,
\end{equation}
y
\begin{equation}
\omega \gg \omega _{n}\rightarrow 20\log _{10}\left\| H\left( j\omega
\right) \right\| \approx -40\log _{10}\omega +40\log _{10}\omega _{n}.
\end{equation}
\end{enumerate}

\section{Diagramas de Bode para respuestas en frecuencias racionales}

Supongamos que deseamos realizar el diagrama de Bode para un sistema cuya
funci\'{o}n de transferencia es 
\begin{equation}
H\left( s\right) =\frac{P\left( s\right) }{Q\left( s\right) }=\frac{%
\sum_{i=0}^{N}a_{i}s^{i}}{\sum_{j=0}^{M}b_{j}s^{j}},
\end{equation}
donde $M>N.$ Para poder hacer el gr\'{a}fico de la magnitud y de la fase de la
respuesta en frecuencia nos conviene llevar a cabo los siguientes pasos:

\begin{enumerate}
\item  Encontrar los ceros de los polin\'{o}mios $P\left( s\right) $ y $Q\left(
s\right) $. Recordemos que un polin\'{o}mio con coeficientes reales tiene ceros
reales o complejos conjugados. Lo que nos interesa es escribir los
polin\'{o}mios como producto de factores con las formas:

\begin{enumerate}
\item  $m,$ donde $m$ es un n\'{u}mero real.

\item  $s,$ cuando la ra\'{\i}z del polin\'{o}mio es nula.

\item  $\left( s+c\right) $ donde $c$ es un n\'{u}mero real no nulo.

\item  $\left( s+d\right) ^{2}+e^{2}$ donde $-d+je$ y $-d-je$
son ra\'{\i}ces complejas conjugadas.
\end{enumerate}

\item  Supongamos que tenemos los polin\'{o}mios $P\left( s\right) $ y $Q\left(
s\right) $ tienen la forma general $G\left( s\right) $ 
\begin{equation}
        G\left( s\right) =
                m\left( \prod_{i=1}^{R}s+c_{i}\right) 
                \left(
                        \prod_{k=1}^{S}s^{2}+2d_{k}s+\left( d_{k}^{2}+e_{k}^{2}\right) 
                \right) s^{T},
\end{equation}

\item  Podemos escribir por lo tanto la funci\'{o}n de transferencia como
una productoria de funciones parciales normalizadas 
\begin{equation}
H\left( s\right) =\prod_{k=1}^{U}H_{k}\left( s\right) ,
\end{equation}
donde las funciones $H_{k}\left( s\right) $ tiene la forma

\begin{enumerate}
\item  $B_{1}\left( s\right) =K,$ donde $K$ es un n\'{u}mero real,

\item  $B_{2}\left( s\right) =\frac{1}{\frac{s}{c_{i}}+1 },$

\item  $B_{3}\left( s\right) =\frac{1}{ \frac{s^{2}}{\left(
d_{k}^{2}+e_{k}^{2}\right) }+\frac{2d_{k}}{\left(
d_{k}^{2}+e_{k}^{2}\right) }s+1 },$

\item  $B_{4}\left( s\right) =\frac{1}{s^{T}},$ donde $T$ es un n\'{u}mero
entero,

\item  $B_{5}\left( s\right) = \frac{s}{c_{i}}+1 ,$

\item  $B_{6}\left( s\right) = \frac{s^{2}}{\left(
d_{k}^{2}+e_{k}^{2}\right) }+\frac{2d_{k}}{\left(
d_{k}^{2}+e_{k}^{2}\right) }s+1 ,$

\item  $B_{7}\left( s\right) =s^{T},$ donde $T$ es un n\'{u}mero entero.
\end{enumerate}

\item  Grafiquemos por separado las magnitudes y fases de cada una de las
funciones $H_{k}\left( j\omega \right) .$ Observe que aunque no hemos
estudiado en particular los casos $B_{5},$ $B_{6}$ y $B_{7},$ estamos en
condiciones de obtener las funciones magnitud y fase de Bode correspodientes
dado que se verifican las siguientes igualdades 
\begin{equation}
\left\| B_{2}\left( j\omega \right) \right\| =\frac{1}{\left\| B_{1}\left(
j\omega \right) \right\| },
\end{equation}
y 
\begin{equation}
\angle B_{2}\left( j\omega \right) =-\angle B_{1}\left( j\omega \right) .
\end{equation}

\item  Dada la especial definici\'{o}n de la funci\'{o}n magnitud de Bode y las
propiedades de los n\'{u}meros complejos, podemos obtener las funciones magnitud
y fase del sistema completo sumando las funciones correspondientes de los
factores $H_{k}\left( j\omega \right) .$
\end{enumerate}

\subsection{Ejemplos}

\begin{figure}
  \centering
  \includegraphics[height=3in]{./graphics/bodeone.eps}
  \caption{Diagrama de bode para un sistema con un polo}
  \label{bodeoneeps}
\end{figure}

\begin{figure}
  \centering
  \includegraphics[height=3in]{./graphics/bodetwo.eps}
  \caption{Diagrama de bode para un sistema con dos polos}
  \label{bodetwoeps}
\end{figure}

